% Step 4 interpretation section
% This file is intended to be included with % Step 4 interpretation section
% This file is intended to be included with % Step 4 interpretation section
% This file is intended to be included with % Step 4 interpretation section
% This file is intended to be included with \input{step4_interpretation.tex}
% Adjust figure labels if your main report uses different names.

\section{Interpretation of Step 4 Results and Plots}
\label{sec:step4_interpretation}

\subsection{Summary of the primary result}
The pre-specified primary analysis did not provide statistically significant evidence for a
front-versus-back dissociation between abstract and concrete questions.
Across the $n=36$ included participant-sessions, the primary dissociation score was positive on
average but small, and the one-sample test against zero was not significant,
$t(35)=1.3547$, $p=0.184185$, with a small effect size ($d=0.2258$).
Accordingly, the main conclusion of Step~4 is that the data show a descriptive trend in the expected
spatial direction, but not sufficiently strong evidence to reject the null hypothesis in the primary test.

\subsection{Interpretation of the planned ROI follow-up tests}
The planned follow-up tests help describe the direction of the pattern, even though they do not alter
the primary inferential conclusion.
In the front ROI, the abstract-versus-concrete contrast was positive but not significant,
$t(35)=1.4117$, $p=0.166864$, Holm-adjusted $p=0.333728$, $d_z=0.2353$.
This suggests a small tendency for abstract questions to elicit a larger front-ROI HbO response than
concrete questions, but the evidence is too weak to support a reliable condition difference.
In the back ROI, the abstract-versus-concrete contrast was essentially null,
$t(35)=-0.0496$, $p=0.960692$, Holm-adjusted $p=0.960692$, $d_z=-0.0083$,
indicating no meaningful separation between the two question types in the posterior ROI.

Taken together, the ROI results suggest that any condition effect is driven mainly by a mild frontal
trend rather than by a robust front-versus-back dissociation.

\subsection{Interpretation of the participant-level four-cell summary}
Figure~\ref{fig:step4_fourcell} is useful for understanding why the primary test did not reach
statistical significance.
At the descriptive level, the front ROI shows a small downward shift from abstract to concrete,
whereas the back ROI shows very similar means across the two conditions.
However, the participant-level grey trajectories reveal substantial inter-individual variability, and
many participants do not follow the same direction of change.
This heterogeneity weakens the group-level dissociation test, even though the average pattern is
slightly more favorable to the hypothesized frontal abstract effect than to the opposite pattern.

Another important point is that the baseline-corrected HbO values are very small in absolute units.
For this reason, the substantive interpretation should focus more on the relative ordering of the four
cells and on the consistency of participant-level trajectories than on the rounded mean values alone.

\subsection{Interpretation of the dissociation-score distribution}
Figure~\ref{fig:step4_dissociation_hist} shows the distribution of the participant-level primary
dissociation scores.
The distribution is centered only slightly above zero and still includes many negative or near-zero
values.
This is consistent with the statistical test: the sample does not cluster strongly enough on the
positive side to provide reliable evidence for a robust front-versus-back condition difference.

At the same time, the histogram does not suggest complete absence of pattern.
Instead, it is more consistent with a weak and noisy positive effect that varies considerably across
participants.
This distinction is important for interpretation: Step~4 does not indicate that the effect is impossible
or absent in principle, only that the present primary summary metric does not capture it strongly
enough to yield statistical significance.

\subsection{Interpretation of the paired ROI plots}
Figure~\ref{fig:step4_roi_pairs} provides the clearest visual explanation of the follow-up tests.
In the front ROI, abstract questions tend to sit slightly above concrete questions for the average
participant, but the participant-specific slopes are mixed and frequently shallow.
This matches the small positive but non-significant front-ROI result.

In the back ROI, the two condition means are nearly overlapping, and the participant-level slopes go
in both directions with no consistent group tendency.
This visual pattern is fully consistent with the near-zero back-ROI test statistic.
Therefore, the paired ROI plots reinforce the interpretation that the current Step~4 effect is best
described as a weak frontal tendency rather than a full anterior-posterior dissociation.

\subsection{Interpretation of the descriptive HbO waveforms}
Figure~\ref{fig:step4_waveforms} adds an important temporal perspective.
In the front ROI, the abstract and concrete waveforms begin similarly near onset and then show a
modest separation in the mid-to-late post-onset interval, with the abstract curve tending to remain
above the concrete curve for part of the response window.
In the back ROI, both conditions follow a broadly similar temporal shape, with much less convincing
visual separation.

This temporal structure suggests that the condition difference may be more localized in time than the
primary Step~4 summary metric assumes.
Because the primary analysis averaged from question onset all the way to answer, using trial windows
of variable duration, a difference that is strongest only in a narrower post-onset interval may have
been diluted when collapsed into a single whole-window mean.
Thus, the waveform plots do not overturn the non-significant primary result, but they do suggest a
plausible explanation for why a small descriptive effect is visible in the figures while the inferential
test remains non-significant.

\subsection{Interpretation of the secondary analyses}
The secondary analyses also support a cautious reading of the results.
The analogous HbR primary-style analysis was not significant ($p=0.838320$), which means that the
main result is not corroborated by a clear complementary HbR pattern.
Likewise, the smallest corrected $p$-value in the exploratory channel-wise HbO analysis was
$0.355121$, so no individual long channel survived multiple-comparison correction.

This does not imply that the exploratory map is uninformative.
Rather, it means that the spatial pattern should be treated as descriptive only.
The map in Figure~\ref{fig:step4_channel_map} suggests mixed positive and negative channel-level
contrasts, including a visually more positive anterior pattern and a more heterogeneous posterior
pattern, but the evidence is not strong enough to support channel-specific inferential claims.

\subsection{Interpretation of the supplementary visualization report}
The supplementary plots provide useful context for the main result.
The onset-locked visualization report included the same $36$ participant-sessions used in Step~4,
with $539$ abstract epochs and $536$ concrete epochs, and no trials were dropped for the fixed epoch
window from $-3.5$~s to $13.2$~s around question onset.
The baseline was defined as $[-2,0]$~s, and the report also provided a single-time comparison at
$9.0$~s.

The sensor-geometry figure confirms that the montage samples only the left lateral portion of the
head and that the short pairs are geometrically distinct from the long cortical channels.
This supports the methodological choice made in Step~4 to treat the analysis as a left-lateral ROI
comparison rather than a whole-brain statement.

The onset-locked topography figures for abstract and concrete questions suggest that the two
conditions share broad haemodynamic features, but they also hint at subtle time-dependent spatial
differences.
The single-time comparison at $9.0$~s is especially informative because it indicates that condition
contrasts may emerge more clearly at specific post-onset moments than in the variable-duration
whole-question averages used in the primary test.
This again supports the interpretation that the present non-significant primary result may reflect a
mismatch between the strongest part of the haemodynamic effect and the broad trial summary used in
Step~4.

\subsection{Overall interpretation}
The overall interpretation of Step~4 should be balanced.
The primary pre-specified test was not statistically significant, and therefore the study cannot claim
reliable inferential evidence for a front-versus-back dissociation between abstract and concrete
questions at this stage.
At the same time, the descriptive plots are not purely flat or random.
They suggest a small frontal tendency in the expected direction and a more time-localized pattern
that may not be captured optimally by the current primary summary metric.

A careful summary is therefore as follows:
\begin{quote}
Step~4 revealed a weak descriptive pattern consistent with greater frontal engagement during abstract
questions, but this pattern did not reach statistical significance in the pre-specified primary
front-versus-back dissociation test.
\end{quote}

\subsection{Recommended wording for the report conclusion}
A concise conclusion paragraph could read as follows:

\begin{quote}
In the pre-specified primary ROI analysis, the front-versus-back dissociation between abstract and
concrete questions did not reach statistical significance, $t(35)=1.3547$, $p=0.184185$,
$d=0.2258$.
Planned follow-up tests showed only a small non-significant tendency for abstract questions to elicit
higher HbO responses than concrete questions in the front ROI, while the back ROI showed virtually
no condition difference.
Secondary HbR and corrected channel-wise analyses were likewise non-significant.
Overall, the Step~4 results support a weak descriptive effect pattern rather than robust inferential
evidence for an abstract-versus-concrete dissociation in the current primary analysis.
\end{quote}

\subsection{Bridge to the next step}
The plots suggest a useful direction for the next analysis stage.
Because the onset-locked waveforms and the single-time topographic comparison indicate that the
condition effect may be temporally localized, a natural next step is to test a fixed post-onset window
or a small number of pre-specified onset-locked windows rather than averaging from question onset to
the participant's response time.
This would still remain relatively simple while better matching the temporal profile suggested by the
supplementary visualizations.

% Adjust figure labels if your main report uses different names.

\section{Interpretation of Step 4 Results and Plots}
\label{sec:step4_interpretation}

\subsection{Summary of the primary result}
The pre-specified primary analysis did not provide statistically significant evidence for a
front-versus-back dissociation between abstract and concrete questions.
Across the $n=36$ included participant-sessions, the primary dissociation score was positive on
average but small, and the one-sample test against zero was not significant,
$t(35)=1.3547$, $p=0.184185$, with a small effect size ($d=0.2258$).
Accordingly, the main conclusion of Step~4 is that the data show a descriptive trend in the expected
spatial direction, but not sufficiently strong evidence to reject the null hypothesis in the primary test.

\subsection{Interpretation of the planned ROI follow-up tests}
The planned follow-up tests help describe the direction of the pattern, even though they do not alter
the primary inferential conclusion.
In the front ROI, the abstract-versus-concrete contrast was positive but not significant,
$t(35)=1.4117$, $p=0.166864$, Holm-adjusted $p=0.333728$, $d_z=0.2353$.
This suggests a small tendency for abstract questions to elicit a larger front-ROI HbO response than
concrete questions, but the evidence is too weak to support a reliable condition difference.
In the back ROI, the abstract-versus-concrete contrast was essentially null,
$t(35)=-0.0496$, $p=0.960692$, Holm-adjusted $p=0.960692$, $d_z=-0.0083$,
indicating no meaningful separation between the two question types in the posterior ROI.

Taken together, the ROI results suggest that any condition effect is driven mainly by a mild frontal
trend rather than by a robust front-versus-back dissociation.

\subsection{Interpretation of the participant-level four-cell summary}
Figure~\ref{fig:step4_fourcell} is useful for understanding why the primary test did not reach
statistical significance.
At the descriptive level, the front ROI shows a small downward shift from abstract to concrete,
whereas the back ROI shows very similar means across the two conditions.
However, the participant-level grey trajectories reveal substantial inter-individual variability, and
many participants do not follow the same direction of change.
This heterogeneity weakens the group-level dissociation test, even though the average pattern is
slightly more favorable to the hypothesized frontal abstract effect than to the opposite pattern.

Another important point is that the baseline-corrected HbO values are very small in absolute units.
For this reason, the substantive interpretation should focus more on the relative ordering of the four
cells and on the consistency of participant-level trajectories than on the rounded mean values alone.

\subsection{Interpretation of the dissociation-score distribution}
Figure~\ref{fig:step4_dissociation_hist} shows the distribution of the participant-level primary
dissociation scores.
The distribution is centered only slightly above zero and still includes many negative or near-zero
values.
This is consistent with the statistical test: the sample does not cluster strongly enough on the
positive side to provide reliable evidence for a robust front-versus-back condition difference.

At the same time, the histogram does not suggest complete absence of pattern.
Instead, it is more consistent with a weak and noisy positive effect that varies considerably across
participants.
This distinction is important for interpretation: Step~4 does not indicate that the effect is impossible
or absent in principle, only that the present primary summary metric does not capture it strongly
enough to yield statistical significance.

\subsection{Interpretation of the paired ROI plots}
Figure~\ref{fig:step4_roi_pairs} provides the clearest visual explanation of the follow-up tests.
In the front ROI, abstract questions tend to sit slightly above concrete questions for the average
participant, but the participant-specific slopes are mixed and frequently shallow.
This matches the small positive but non-significant front-ROI result.

In the back ROI, the two condition means are nearly overlapping, and the participant-level slopes go
in both directions with no consistent group tendency.
This visual pattern is fully consistent with the near-zero back-ROI test statistic.
Therefore, the paired ROI plots reinforce the interpretation that the current Step~4 effect is best
described as a weak frontal tendency rather than a full anterior-posterior dissociation.

\subsection{Interpretation of the descriptive HbO waveforms}
Figure~\ref{fig:step4_waveforms} adds an important temporal perspective.
In the front ROI, the abstract and concrete waveforms begin similarly near onset and then show a
modest separation in the mid-to-late post-onset interval, with the abstract curve tending to remain
above the concrete curve for part of the response window.
In the back ROI, both conditions follow a broadly similar temporal shape, with much less convincing
visual separation.

This temporal structure suggests that the condition difference may be more localized in time than the
primary Step~4 summary metric assumes.
Because the primary analysis averaged from question onset all the way to answer, using trial windows
of variable duration, a difference that is strongest only in a narrower post-onset interval may have
been diluted when collapsed into a single whole-window mean.
Thus, the waveform plots do not overturn the non-significant primary result, but they do suggest a
plausible explanation for why a small descriptive effect is visible in the figures while the inferential
test remains non-significant.

\subsection{Interpretation of the secondary analyses}
The secondary analyses also support a cautious reading of the results.
The analogous HbR primary-style analysis was not significant ($p=0.838320$), which means that the
main result is not corroborated by a clear complementary HbR pattern.
Likewise, the smallest corrected $p$-value in the exploratory channel-wise HbO analysis was
$0.355121$, so no individual long channel survived multiple-comparison correction.

This does not imply that the exploratory map is uninformative.
Rather, it means that the spatial pattern should be treated as descriptive only.
The map in Figure~\ref{fig:step4_channel_map} suggests mixed positive and negative channel-level
contrasts, including a visually more positive anterior pattern and a more heterogeneous posterior
pattern, but the evidence is not strong enough to support channel-specific inferential claims.

\subsection{Interpretation of the supplementary visualization report}
The supplementary plots provide useful context for the main result.
The onset-locked visualization report included the same $36$ participant-sessions used in Step~4,
with $539$ abstract epochs and $536$ concrete epochs, and no trials were dropped for the fixed epoch
window from $-3.5$~s to $13.2$~s around question onset.
The baseline was defined as $[-2,0]$~s, and the report also provided a single-time comparison at
$9.0$~s.

The sensor-geometry figure confirms that the montage samples only the left lateral portion of the
head and that the short pairs are geometrically distinct from the long cortical channels.
This supports the methodological choice made in Step~4 to treat the analysis as a left-lateral ROI
comparison rather than a whole-brain statement.

The onset-locked topography figures for abstract and concrete questions suggest that the two
conditions share broad haemodynamic features, but they also hint at subtle time-dependent spatial
differences.
The single-time comparison at $9.0$~s is especially informative because it indicates that condition
contrasts may emerge more clearly at specific post-onset moments than in the variable-duration
whole-question averages used in the primary test.
This again supports the interpretation that the present non-significant primary result may reflect a
mismatch between the strongest part of the haemodynamic effect and the broad trial summary used in
Step~4.

\subsection{Overall interpretation}
The overall interpretation of Step~4 should be balanced.
The primary pre-specified test was not statistically significant, and therefore the study cannot claim
reliable inferential evidence for a front-versus-back dissociation between abstract and concrete
questions at this stage.
At the same time, the descriptive plots are not purely flat or random.
They suggest a small frontal tendency in the expected direction and a more time-localized pattern
that may not be captured optimally by the current primary summary metric.

A careful summary is therefore as follows:
\begin{quote}
Step~4 revealed a weak descriptive pattern consistent with greater frontal engagement during abstract
questions, but this pattern did not reach statistical significance in the pre-specified primary
front-versus-back dissociation test.
\end{quote}

\subsection{Recommended wording for the report conclusion}
A concise conclusion paragraph could read as follows:

\begin{quote}
In the pre-specified primary ROI analysis, the front-versus-back dissociation between abstract and
concrete questions did not reach statistical significance, $t(35)=1.3547$, $p=0.184185$,
$d=0.2258$.
Planned follow-up tests showed only a small non-significant tendency for abstract questions to elicit
higher HbO responses than concrete questions in the front ROI, while the back ROI showed virtually
no condition difference.
Secondary HbR and corrected channel-wise analyses were likewise non-significant.
Overall, the Step~4 results support a weak descriptive effect pattern rather than robust inferential
evidence for an abstract-versus-concrete dissociation in the current primary analysis.
\end{quote}

\subsection{Bridge to the next step}
The plots suggest a useful direction for the next analysis stage.
Because the onset-locked waveforms and the single-time topographic comparison indicate that the
condition effect may be temporally localized, a natural next step is to test a fixed post-onset window
or a small number of pre-specified onset-locked windows rather than averaging from question onset to
the participant's response time.
This would still remain relatively simple while better matching the temporal profile suggested by the
supplementary visualizations.

% Adjust figure labels if your main report uses different names.

\section{Interpretation of Step 4 Results and Plots}
\label{sec:step4_interpretation}

\subsection{Summary of the primary result}
The pre-specified primary analysis did not provide statistically significant evidence for a
front-versus-back dissociation between abstract and concrete questions.
Across the $n=36$ included participant-sessions, the primary dissociation score was positive on
average but small, and the one-sample test against zero was not significant,
$t(35)=1.3547$, $p=0.184185$, with a small effect size ($d=0.2258$).
Accordingly, the main conclusion of Step~4 is that the data show a descriptive trend in the expected
spatial direction, but not sufficiently strong evidence to reject the null hypothesis in the primary test.

\subsection{Interpretation of the planned ROI follow-up tests}
The planned follow-up tests help describe the direction of the pattern, even though they do not alter
the primary inferential conclusion.
In the front ROI, the abstract-versus-concrete contrast was positive but not significant,
$t(35)=1.4117$, $p=0.166864$, Holm-adjusted $p=0.333728$, $d_z=0.2353$.
This suggests a small tendency for abstract questions to elicit a larger front-ROI HbO response than
concrete questions, but the evidence is too weak to support a reliable condition difference.
In the back ROI, the abstract-versus-concrete contrast was essentially null,
$t(35)=-0.0496$, $p=0.960692$, Holm-adjusted $p=0.960692$, $d_z=-0.0083$,
indicating no meaningful separation between the two question types in the posterior ROI.

Taken together, the ROI results suggest that any condition effect is driven mainly by a mild frontal
trend rather than by a robust front-versus-back dissociation.

\subsection{Interpretation of the participant-level four-cell summary}
Figure~\ref{fig:step4_fourcell} is useful for understanding why the primary test did not reach
statistical significance.
At the descriptive level, the front ROI shows a small downward shift from abstract to concrete,
whereas the back ROI shows very similar means across the two conditions.
However, the participant-level grey trajectories reveal substantial inter-individual variability, and
many participants do not follow the same direction of change.
This heterogeneity weakens the group-level dissociation test, even though the average pattern is
slightly more favorable to the hypothesized frontal abstract effect than to the opposite pattern.

Another important point is that the baseline-corrected HbO values are very small in absolute units.
For this reason, the substantive interpretation should focus more on the relative ordering of the four
cells and on the consistency of participant-level trajectories than on the rounded mean values alone.

\subsection{Interpretation of the dissociation-score distribution}
Figure~\ref{fig:step4_dissociation_hist} shows the distribution of the participant-level primary
dissociation scores.
The distribution is centered only slightly above zero and still includes many negative or near-zero
values.
This is consistent with the statistical test: the sample does not cluster strongly enough on the
positive side to provide reliable evidence for a robust front-versus-back condition difference.

At the same time, the histogram does not suggest complete absence of pattern.
Instead, it is more consistent with a weak and noisy positive effect that varies considerably across
participants.
This distinction is important for interpretation: Step~4 does not indicate that the effect is impossible
or absent in principle, only that the present primary summary metric does not capture it strongly
enough to yield statistical significance.

\subsection{Interpretation of the paired ROI plots}
Figure~\ref{fig:step4_roi_pairs} provides the clearest visual explanation of the follow-up tests.
In the front ROI, abstract questions tend to sit slightly above concrete questions for the average
participant, but the participant-specific slopes are mixed and frequently shallow.
This matches the small positive but non-significant front-ROI result.

In the back ROI, the two condition means are nearly overlapping, and the participant-level slopes go
in both directions with no consistent group tendency.
This visual pattern is fully consistent with the near-zero back-ROI test statistic.
Therefore, the paired ROI plots reinforce the interpretation that the current Step~4 effect is best
described as a weak frontal tendency rather than a full anterior-posterior dissociation.

\subsection{Interpretation of the descriptive HbO waveforms}
Figure~\ref{fig:step4_waveforms} adds an important temporal perspective.
In the front ROI, the abstract and concrete waveforms begin similarly near onset and then show a
modest separation in the mid-to-late post-onset interval, with the abstract curve tending to remain
above the concrete curve for part of the response window.
In the back ROI, both conditions follow a broadly similar temporal shape, with much less convincing
visual separation.

This temporal structure suggests that the condition difference may be more localized in time than the
primary Step~4 summary metric assumes.
Because the primary analysis averaged from question onset all the way to answer, using trial windows
of variable duration, a difference that is strongest only in a narrower post-onset interval may have
been diluted when collapsed into a single whole-window mean.
Thus, the waveform plots do not overturn the non-significant primary result, but they do suggest a
plausible explanation for why a small descriptive effect is visible in the figures while the inferential
test remains non-significant.

\subsection{Interpretation of the secondary analyses}
The secondary analyses also support a cautious reading of the results.
The analogous HbR primary-style analysis was not significant ($p=0.838320$), which means that the
main result is not corroborated by a clear complementary HbR pattern.
Likewise, the smallest corrected $p$-value in the exploratory channel-wise HbO analysis was
$0.355121$, so no individual long channel survived multiple-comparison correction.

This does not imply that the exploratory map is uninformative.
Rather, it means that the spatial pattern should be treated as descriptive only.
The map in Figure~\ref{fig:step4_channel_map} suggests mixed positive and negative channel-level
contrasts, including a visually more positive anterior pattern and a more heterogeneous posterior
pattern, but the evidence is not strong enough to support channel-specific inferential claims.

\subsection{Interpretation of the supplementary visualization report}
The supplementary plots provide useful context for the main result.
The onset-locked visualization report included the same $36$ participant-sessions used in Step~4,
with $539$ abstract epochs and $536$ concrete epochs, and no trials were dropped for the fixed epoch
window from $-3.5$~s to $13.2$~s around question onset.
The baseline was defined as $[-2,0]$~s, and the report also provided a single-time comparison at
$9.0$~s.

The sensor-geometry figure confirms that the montage samples only the left lateral portion of the
head and that the short pairs are geometrically distinct from the long cortical channels.
This supports the methodological choice made in Step~4 to treat the analysis as a left-lateral ROI
comparison rather than a whole-brain statement.

The onset-locked topography figures for abstract and concrete questions suggest that the two
conditions share broad haemodynamic features, but they also hint at subtle time-dependent spatial
differences.
The single-time comparison at $9.0$~s is especially informative because it indicates that condition
contrasts may emerge more clearly at specific post-onset moments than in the variable-duration
whole-question averages used in the primary test.
This again supports the interpretation that the present non-significant primary result may reflect a
mismatch between the strongest part of the haemodynamic effect and the broad trial summary used in
Step~4.

\subsection{Overall interpretation}
The overall interpretation of Step~4 should be balanced.
The primary pre-specified test was not statistically significant, and therefore the study cannot claim
reliable inferential evidence for a front-versus-back dissociation between abstract and concrete
questions at this stage.
At the same time, the descriptive plots are not purely flat or random.
They suggest a small frontal tendency in the expected direction and a more time-localized pattern
that may not be captured optimally by the current primary summary metric.

A careful summary is therefore as follows:
\begin{quote}
Step~4 revealed a weak descriptive pattern consistent with greater frontal engagement during abstract
questions, but this pattern did not reach statistical significance in the pre-specified primary
front-versus-back dissociation test.
\end{quote}

\subsection{Recommended wording for the report conclusion}
A concise conclusion paragraph could read as follows:

\begin{quote}
In the pre-specified primary ROI analysis, the front-versus-back dissociation between abstract and
concrete questions did not reach statistical significance, $t(35)=1.3547$, $p=0.184185$,
$d=0.2258$.
Planned follow-up tests showed only a small non-significant tendency for abstract questions to elicit
higher HbO responses than concrete questions in the front ROI, while the back ROI showed virtually
no condition difference.
Secondary HbR and corrected channel-wise analyses were likewise non-significant.
Overall, the Step~4 results support a weak descriptive effect pattern rather than robust inferential
evidence for an abstract-versus-concrete dissociation in the current primary analysis.
\end{quote}

\subsection{Bridge to the next step}
The plots suggest a useful direction for the next analysis stage.
Because the onset-locked waveforms and the single-time topographic comparison indicate that the
condition effect may be temporally localized, a natural next step is to test a fixed post-onset window
or a small number of pre-specified onset-locked windows rather than averaging from question onset to
the participant's response time.
This would still remain relatively simple while better matching the temporal profile suggested by the
supplementary visualizations.

% Adjust figure labels if your main report uses different names.

\section{Interpretation of Step 4 Results and Plots}
\label{sec:step4_interpretation}

\subsection{Summary of the primary result}
The pre-specified primary analysis did not provide statistically significant evidence for a
front-versus-back dissociation between abstract and concrete questions.
Across the $n=36$ included participant-sessions, the primary dissociation score was positive on
average but small, and the one-sample test against zero was not significant,
$t(35)=1.3547$, $p=0.184185$, with a small effect size ($d=0.2258$).
Accordingly, the main conclusion of Step~4 is that the data show a descriptive trend in the expected
spatial direction, but not sufficiently strong evidence to reject the null hypothesis in the primary test.

\subsection{Interpretation of the planned ROI follow-up tests}
The planned follow-up tests help describe the direction of the pattern, even though they do not alter
the primary inferential conclusion.
In the front ROI, the abstract-versus-concrete contrast was positive but not significant,
$t(35)=1.4117$, $p=0.166864$, Holm-adjusted $p=0.333728$, $d_z=0.2353$.
This suggests a small tendency for abstract questions to elicit a larger front-ROI HbO response than
concrete questions, but the evidence is too weak to support a reliable condition difference.
In the back ROI, the abstract-versus-concrete contrast was essentially null,
$t(35)=-0.0496$, $p=0.960692$, Holm-adjusted $p=0.960692$, $d_z=-0.0083$,
indicating no meaningful separation between the two question types in the posterior ROI.

Taken together, the ROI results suggest that any condition effect is driven mainly by a mild frontal
trend rather than by a robust front-versus-back dissociation.

\subsection{Interpretation of the participant-level four-cell summary}
Figure~\ref{fig:step4_fourcell} is useful for understanding why the primary test did not reach
statistical significance.
At the descriptive level, the front ROI shows a small downward shift from abstract to concrete,
whereas the back ROI shows very similar means across the two conditions.
However, the participant-level grey trajectories reveal substantial inter-individual variability, and
many participants do not follow the same direction of change.
This heterogeneity weakens the group-level dissociation test, even though the average pattern is
slightly more favorable to the hypothesized frontal abstract effect than to the opposite pattern.

Another important point is that the baseline-corrected HbO values are very small in absolute units.
For this reason, the substantive interpretation should focus more on the relative ordering of the four
cells and on the consistency of participant-level trajectories than on the rounded mean values alone.

\subsection{Interpretation of the dissociation-score distribution}
Figure~\ref{fig:step4_dissociation_hist} shows the distribution of the participant-level primary
dissociation scores.
The distribution is centered only slightly above zero and still includes many negative or near-zero
values.
This is consistent with the statistical test: the sample does not cluster strongly enough on the
positive side to provide reliable evidence for a robust front-versus-back condition difference.

At the same time, the histogram does not suggest complete absence of pattern.
Instead, it is more consistent with a weak and noisy positive effect that varies considerably across
participants.
This distinction is important for interpretation: Step~4 does not indicate that the effect is impossible
or absent in principle, only that the present primary summary metric does not capture it strongly
enough to yield statistical significance.

\subsection{Interpretation of the paired ROI plots}
Figure~\ref{fig:step4_roi_pairs} provides the clearest visual explanation of the follow-up tests.
In the front ROI, abstract questions tend to sit slightly above concrete questions for the average
participant, but the participant-specific slopes are mixed and frequently shallow.
This matches the small positive but non-significant front-ROI result.

In the back ROI, the two condition means are nearly overlapping, and the participant-level slopes go
in both directions with no consistent group tendency.
This visual pattern is fully consistent with the near-zero back-ROI test statistic.
Therefore, the paired ROI plots reinforce the interpretation that the current Step~4 effect is best
described as a weak frontal tendency rather than a full anterior-posterior dissociation.

\subsection{Interpretation of the descriptive HbO waveforms}
Figure~\ref{fig:step4_waveforms} adds an important temporal perspective.
In the front ROI, the abstract and concrete waveforms begin similarly near onset and then show a
modest separation in the mid-to-late post-onset interval, with the abstract curve tending to remain
above the concrete curve for part of the response window.
In the back ROI, both conditions follow a broadly similar temporal shape, with much less convincing
visual separation.

This temporal structure suggests that the condition difference may be more localized in time than the
primary Step~4 summary metric assumes.
Because the primary analysis averaged from question onset all the way to answer, using trial windows
of variable duration, a difference that is strongest only in a narrower post-onset interval may have
been diluted when collapsed into a single whole-window mean.
Thus, the waveform plots do not overturn the non-significant primary result, but they do suggest a
plausible explanation for why a small descriptive effect is visible in the figures while the inferential
test remains non-significant.

\subsection{Interpretation of the secondary analyses}
The secondary analyses also support a cautious reading of the results.
The analogous HbR primary-style analysis was not significant ($p=0.838320$), which means that the
main result is not corroborated by a clear complementary HbR pattern.
Likewise, the smallest corrected $p$-value in the exploratory channel-wise HbO analysis was
$0.355121$, so no individual long channel survived multiple-comparison correction.

This does not imply that the exploratory map is uninformative.
Rather, it means that the spatial pattern should be treated as descriptive only.
The map in Figure~\ref{fig:step4_channel_map} suggests mixed positive and negative channel-level
contrasts, including a visually more positive anterior pattern and a more heterogeneous posterior
pattern, but the evidence is not strong enough to support channel-specific inferential claims.

\subsection{Interpretation of the supplementary visualization report}
The supplementary plots provide useful context for the main result.
The onset-locked visualization report included the same $36$ participant-sessions used in Step~4,
with $539$ abstract epochs and $536$ concrete epochs, and no trials were dropped for the fixed epoch
window from $-3.5$~s to $13.2$~s around question onset.
The baseline was defined as $[-2,0]$~s, and the report also provided a single-time comparison at
$9.0$~s.

The sensor-geometry figure confirms that the montage samples only the left lateral portion of the
head and that the short pairs are geometrically distinct from the long cortical channels.
This supports the methodological choice made in Step~4 to treat the analysis as a left-lateral ROI
comparison rather than a whole-brain statement.

The onset-locked topography figures for abstract and concrete questions suggest that the two
conditions share broad haemodynamic features, but they also hint at subtle time-dependent spatial
differences.
The single-time comparison at $9.0$~s is especially informative because it indicates that condition
contrasts may emerge more clearly at specific post-onset moments than in the variable-duration
whole-question averages used in the primary test.
This again supports the interpretation that the present non-significant primary result may reflect a
mismatch between the strongest part of the haemodynamic effect and the broad trial summary used in
Step~4.

\subsection{Overall interpretation}
The overall interpretation of Step~4 should be balanced.
The primary pre-specified test was not statistically significant, and therefore the study cannot claim
reliable inferential evidence for a front-versus-back dissociation between abstract and concrete
questions at this stage.
At the same time, the descriptive plots are not purely flat or random.
They suggest a small frontal tendency in the expected direction and a more time-localized pattern
that may not be captured optimally by the current primary summary metric.

A careful summary is therefore as follows:
\begin{quote}
Step~4 revealed a weak descriptive pattern consistent with greater frontal engagement during abstract
questions, but this pattern did not reach statistical significance in the pre-specified primary
front-versus-back dissociation test.
\end{quote}

\subsection{Recommended wording for the report conclusion}
A concise conclusion paragraph could read as follows:

\begin{quote}
In the pre-specified primary ROI analysis, the front-versus-back dissociation between abstract and
concrete questions did not reach statistical significance, $t(35)=1.3547$, $p=0.184185$,
$d=0.2258$.
Planned follow-up tests showed only a small non-significant tendency for abstract questions to elicit
higher HbO responses than concrete questions in the front ROI, while the back ROI showed virtually
no condition difference.
Secondary HbR and corrected channel-wise analyses were likewise non-significant.
Overall, the Step~4 results support a weak descriptive effect pattern rather than robust inferential
evidence for an abstract-versus-concrete dissociation in the current primary analysis.
\end{quote}

\subsection{Bridge to the next step}
The plots suggest a useful direction for the next analysis stage.
Because the onset-locked waveforms and the single-time topographic comparison indicate that the
condition effect may be temporally localized, a natural next step is to test a fixed post-onset window
or a small number of pre-specified onset-locked windows rather than averaging from question onset to
the participant's response time.
This would still remain relatively simple while better matching the temporal profile suggested by the
supplementary visualizations.
